\section{Searching}
\paragraph{Binary Search}
\textcolor{Blue}{Precondition}: Array is sorted, and of size $n$\\
\textcolor{Blue}{Postcondition}: If element is in arr, \verb|A[begin]=key|\\
\textcolor{Blue}{Invariants}: Key is within \verb|A[begin..end]|\\
\verb|(end-begin)|$\le \frac{n}{2^k}$ after the $k$-th iteration
\begin{verbatim}
int search(A, key, n)
   begin = 0
   end = n-1
   while begin < end do:
      mid = begin + (end - begin)/2
      if key <= A[mid] then
         end = mid
      else begin = mid + 1
   return (A[begin]==key) ? begin : -1
\end{verbatim}
\textcolor{Bittersweet}{Runtime}: $O(\log n)$

\paragraph{Peakfinding}
\textcolor{Blue}{Precondition}: Array of size $n$\\
\textcolor{Blue}{Postcondition}: $A[i-1]\le A[i]\wedge A[i+1]\le A[i]$\\
\textcolor{Blue}{Invariants}: $\exists$ a peak in \verb|A[begin..end]|\\
Every peak in \verb|A[begin..end]| is a peak in \verb|A[0..n-1]|\\
Recurse in the right$\Rightarrow$Peak in the right half is a peak in the array.
\begin{verbatim}
int findPeak(A, n)
   if A[n/2] is a peak then return n/2
   else if A[n/2 + 1] > A[n/2] then
      return findPeak(A[n/2 + 1..n], n/2)
   else if A[n/2 - 1] > A[n/2] then
      return findPeak(A[1..n/2 - 1], n/2)
\end{verbatim}
\textcolor{Bittersweet}{Runtime}: $O(\log n)$\ \textcolor{Bittersweet}{Steep peakfinding}: $O(n)$

\paragraph{Quickselect}
\textcolor{Blue}{Precondition}: Array of size $n$
\textcolor{Blue}{Postcondition}: Return the k-th smallest element.
\textcolor{Blue}{Invariants}: The k-th element is on the side that we recurse on.\\
\textcolor{Bittersweet}{Runtime}: $O(n)$
